\section{How this paper was checked}\label{app:verification}

\paragraph{Sources.}
\begin{itemize}
\item The repository \texttt{KokunoYumeto/erdos-straus-foundation} at commit \texttt{d20b32e}, read-only.
\item The 619-page archive and its source archive, downloaded on 26 September 2026 from the Zenodo record 10.5281/zenodo.22678971:\\[2pt]
{\footnotesize\begin{tabular}{@{}ll@{}}
\texttt{00\_ERDOS\_STRAUSS\_Project\_Reader.pdf} & MD5 \texttt{c1ee553967e0e6f12d4976e171c11b15}\\
\texttt{01\_ERDOS\_STRAUSS\_Source\_Verification\_and\_Machine\_Memory.zip} & MD5 \texttt{4129ec50d0c02ff7daf61fde6c589de4}
\end{tabular}}\\[2pt]
The main source file of the archive (60{,}561 lines) and its fourteen satellite files were read where they are cited.
\item The supplement's source, in the repository folder\newline \texttt{research/expanded-2026-09-08/exact\_bounded\_transport\_72/}.
\item For the literature statements of §1: the source of Elsholtz and Tao's paper (arXiv 1107.1010v6), which the repository holds in its literature folders.
\end{itemize}

\paragraph{What was refereed earlier.}
Theorem~\ref{thm:shell} and Propositions~\ref{prop:r3}--\ref{prop:mod840}, the audited items of Section~\ref{sec:items}, and the audit of the package ES-09 in Section~\ref{ssec:es-s6} come from the Erdős--Straus section of the reader of four workbenches. There they went through five stages:
\begin{enumerate}[label=\arabic*.]
\item a first-pass audit by a separate Claude instance, which read the workbench and wrote independent checks (\note{43a\_});
\item a referee pass on the notes \note{43\_}--\note{47\_}, by an instance that had written none of them; it found, among other things, that the orbit statements of Section~\ref{ssec:es-s6} fail at even levels;
\item a referee pass on that reader;
\item two further passes that read it for understatement and missed generality as well as for overstatement.
\end{enumerate}
These passes produced the range $x<p/2$ in Theorem~\ref{thm:shell}, the first form of Proposition~\ref{prop:survivors}, Proposition~\ref{prop:mod840}, and the extension of the count of Proposition~\ref{prop:r3} to $p\equiv5\pmod{12}$. The checks are \note{checks/\allowbreak workbenches/\allowbreak workbench\_reader\_claims\_checks.py} and \note{checks/\allowbreak workbenches/\allowbreak generality\_checks.py}; the referee reports are in \note{checks/workbenches/}.

\paragraph{What version~1 added to the reader of four workbenches.}
\emph{New} here means new relative to that earlier reader, not a claim of new mathematics. Where the archive already has a statement, this paper cites it.
\begin{itemize}
\item Proposition~\ref{prop:jacobi}, which is the archive's Theorem~10.9 and Proposition~10.10, re-proved;
\item Lemma~\ref{lem:leastclass} and Proposition~\ref{prop:records}, with the reproduction of the deficit data in §\ref{ssec:records};
\item Theorem~\ref{thm:families}, and the statement that each family meets every hard class;
\item Proposition~\ref{prop:twoshears};
\item Proposition~\ref{prop:quartic}(a), whose relation was also checked in \note{21\_};
\item the map of the archive (§\ref{ssec:archive}), the consistency of the census counts (§\ref{ssec:census}), and the three new rows of the overlap table (Section~\ref{sec:overlaps}).
\end{itemize}
These went through the referee pass described next.

\paragraph{The referee pass of version~1.}
An independent Claude instance (\texttt{claude-opus-5-5}), which wrote none of the text, refereed version~1 on 27 September 2026. It read in both directions, as the author asked: for statements stronger than their proofs or sources, and for understatement, missed generality and missed or implied results. It reran the five scripts of version~1, whose outputs were identical, and wrote its own checks. It found every proof correct as written and every recomputed number in agreement. It reported five major and ten minor findings, and seven results that the reader's material implies. Every finding was verified here before it was applied: by reading the archive at the cited page, or by an independent script.
\begin{itemize}
\item \emph{Attributions.} Theorem~\ref{thm:shell}(b) is not a generalisation of the archive. The archive's Theorems~9.1 and~9.64 already hold for every odd prime; the paper now says so. The check here also found that the count $2|E_a|+|M_a|$ is in the archive's Theorem~9.64, which the referee had credited to version~1.
\item \emph{Understatements.} Proposition~\ref{prop:survivors} now concludes that the survivors lie in the hard classes modulo $840$, as the archive's Corollary~10.7 does. The census domain is stated. The page map of §24 is split, and the page counts are made consistent.
\item \emph{Two misdescriptions.}
\begin{itemize}
\item A remark of version~1, that the census carriers are not polynomial identities alone, was wrong. They are polynomial identities on their rows, and Proposition~\ref{prop:nosquare} and Corollary~\ref{cor:censuslimit} replace it.
\item Proposition~\ref{prop:twoshears}(b) had suggested that the theorem fails for every $p\equiv2\pmod3$. It fails only on a minority of them, and Proposition~\ref{prop:sheargraph} and Remark~\ref{rem:sheardensity} now describe that set.
\end{itemize}
\item \emph{Minor findings.} Among them:
\begin{itemize}
\item the residual $R=2$ at $p=2$;
\item the archive's correction of Ionascu and Wilson's label for $1201$;
\item the count of \texttt{.lean} files ($67$, not $84$);
\item the title-page date ($30$ August);
\item Elsholtz and Tao's definition of solvability by polynomials;
\item the full citation of Vaughan's paper.
\end{itemize}
\item \emph{Results adopted.} After verification, these are Corollary~\ref{cor:hle2}, Propositions~\ref{prop:groupbarrier}, \ref{prop:groupconverse}, \ref{prop:sheargraph}, \ref{prop:fixeddivisor} and~\ref{prop:nosquare}, Corollary~\ref{cor:censuslimit}, Proposition~\ref{prop:quartic}(b), the records to $3\cdot10^9$, and the even levels of §\ref{ssec:es-s6}.
\item \emph{A correction to the referee.} It gave $3{,}533{,}141$ as the prime with a shear path of length $5$. A complete enumeration here shows that the least such prime is $1{,}183{,}451$.
\item \emph{Not adopted.} The referee also reported an observation of Ionascu and Wilson on residues modulo $9240$. It could not be checked against their paper here, and is omitted.
\end{itemize}
The referee's report and its check programs are in \note{checks/es\_reader/referee/}; the verification scripts of version~1 are in \note{checks/es\_reader/referee\_verify/}.

\paragraph{The scripts.}
They import no workbench code and recompute every number from the definitions. They are republished in \texttt{contrib/claude-ab/rewrite/checks/es/}; the folders \note{referee/} and \note{referee\_verify/} are in the branch history (see below).

\begin{center}\small
\begin{tabularx}{\linewidth}{@{}>{\raggedright\arraybackslash}p{0.25\linewidth}>{\raggedright\arraybackslash}X>{\raggedright\arraybackslash}p{0.13\linewidth}@{}}
\toprule
Script & What it checks & Time, result\\
\midrule
\note{esr\_records.py} & $h(p)$ for every prime up to $8{,}803{,}369$ through Lemma~\ref{lem:leastclass} and Theorem~\ref{thm:shell}; the strict records; $h(806521)=h(118801)=15$ (for $p=2$ the script prints $R=4h-1=3$; the text's $R=2$ is $4x-p$ with the least denominator $x=1$) & 3 s; as stated\\
\note{esr\_lemma\_checks.py} & Proposition~\ref{prop:jacobi} on all $47{,}210$ shells meeting its hypothesis at odd primes below $3000$ (none occupied), and the failure of its converse; Lemma~\ref{lem:leastclass} against a brute-force enumeration of all solutions at the odd primes below $700$ & 25 s; all pass\\
\note{esr\_deficits.py} & the deficit data of §\ref{ssec:records} at $73$, $1129$, $1201$, $21169$, $118{,}801$, $8{,}803{,}369$ and $806{,}521$, their independence of the primitive root, the corollary of archive §16, and the counts $|E_a|=0$, $|M_a|=2$ at $p_*$ (22 checks) & 1 s; all pass\\
\note{esr\_families.py} & the four identities of Theorem~\ref{thm:families} as polynomial identities, exact fractions for $k\le200$, the hard classes met and the least prime in each & 1 s; all pass\\
\note{esr\_twoshears.py} & Proposition~\ref{prop:twoshears}: the full shear graph of every prime below $2000$, and a complete search for two-edge paths up to $10^7$ & 23 s; all pass\\
\note{referee\_verify/\allowbreak shears\_independent.py} & Proposition~\ref{prop:sheargraph}: degrees on the full graph below $1200$; complete enumeration of paths of length $2$ to $5$ below $10^7$; the class $131\bmod180$ & 21 s; as stated\\
\note{referee\_verify/\allowbreak shears\_density\_\allowbreak bound.py} & the bound $\sum_R\delta_R<0.335$ of Remark~\ref{rem:sheardensity} & 5 s; as stated\\
\note{referee\_verify/\allowbreak hsegment.c}, \note{hsegment\_A.c} & $h(p)$ for all primes below $3\cdot10^9$ by a segmented sieve; the records, the primes with $h\ge19$, and Corollary~\ref{cor:hle2} & 2 min; as stated\\
\note{esr\_referee\_checks.py} & Propositions~\ref{prop:survivors}, \ref{prop:groupbarrier}, \ref{prop:groupconverse}; the census domain; the base rows and the square-row count of Corollary~\ref{cor:censuslimit}; the shell at $p=12289$; the even levels of §\ref{ssec:es-s6} & 2 min; all pass\\
\bottomrule
\end{tabularx}
\end{center}

\paragraph{A partial second pass.}
A second-pass audit of the supplement was started with a separate Claude instance and stopped before it wrote its report. Its scripts and outputs are kept in \note{checks/es\_reader/second\_pass\_partial/}, and §\ref{ssec:supplement} reports what they show. The one reported mismatch was traced to the script's own expectation (§\ref{ssec:supplement}).

\paragraph{Version~3: the author's notes.}
Section~\ref{sec:notes} was added in version~3, on 27 September 2026.
\begin{itemize}
\item Six Claude instances, each of which wrote none of this paper, read the notes in full in separate passes. Each tabulated every numbered statement and checked it with its own programs.
\item claude-ab then re-derived every result proved in Section~\ref{sec:notes} and re-checked it with its own programs, which are in \note{checks/es50/}:
\begin{itemize}
\item \note{es50\_checks.py}: 25 items, all pass, in 49 s. Every certificate is turned into $(x,y,z)$ and tested by $4xyz=p(xy+yz+zx)$.
\item \note{es50\_support\_core.py}: an independent reproduction of the counts $4951$ and $2970$ of Theorem~\ref{thm:termcore}(a), in 4 s.
\item \note{es50\_survivors.py}: Table~\ref{tab:survivors}.
\item \note{salez/}: Theorem~\ref{thm:termcore}(b). The enumerator of Salez's seven modular equations was written by one of the reading passes. It gives $3520$ classes at $840\cdot11\cdot19\cdot23$, and $147{,}348$ at Salez's modulus, which is the number he reports.
\end{itemize}
\item The full record, with the passes' findings and the attributions, is note \note{50\_}.
\item \emph{The referee of version 3.} An independent Claude instance, which wrote none of the text, refereed note \note{50\_} and Section~\ref{sec:notes} on 27 September 2026.
\begin{itemize}
\item It re-ran every program, including the author's own scripts, and recomputed the main numbers by independent code. It found every proof correct.
\item It reported three major findings, all applied: one error charged to the wrong note, an open question false as stated (the prime $409$), and this section's own verification status, which was stated too broadly.
\item It reported fifteen minor findings, applied after checking. Among them: Salez's scope, attributions to the notes and the archive, and the reduced form of Theorem~\ref{thm:visiblebox}.
\item It found results that the material implies. The eight further shifts after Table~\ref{tab:survivors} and the equivalence in Proposition~\ref{prop:groupbarrier}(c) are adopted.
\end{itemize}
Its report and programs are in \note{checks/es50/referee/}.
\end{itemize}

\paragraph{Version~4: the new results.}
The results new in version~4 were proved by claude-ab and checked by programs written for them, which are in \texttt{contrib/claude-ab/rewrite/checks/es/}:
\begin{itemize}
\item \note{exact\_shadow\_reps.py}: Theorem~\ref{thm:m27} and Proposition~\ref{prop:fibonacci}, exactly in $\Q(\zeta_{168})$ and $\Q(\zeta_{120})$, together with the agreement of the $S$-matrix formula with a direct evaluation of the series (to $29$ and $30$ digits);
\item \note{m7\_vs\_M27.py} and \note{m5\_vs\_M25.py}: the same $S$-matrices reproduced numerically from the $q$-series to $38$ digits, the search over signed permutations, conjugation and twists, and the failure of $(ST)^3=S^2$ for the pair of the note's cubic transport;
\item \note{misc\_checks.py}: the values of $A_{11}(pa)$ on the primes $p\equiv1\pmod{840}$, the involutions of $GL_2(3)$, and the coordinate formulas of Proposition~\ref{prop:flowboost};
\item \note{es09\_even\_levels.py} and the scripts of the companion $S^6$ paper \cite{SixSpherePaper} (in \texttt{checks/s6/}): Theorem~\ref{thm:orbits} and Corollary~\ref{cor:orbits}.
\end{itemize}
Proposition~\ref{prop:densityzero}, Proposition~\ref{prop:dossier} and Proposition~\ref{prop:torsor} are proved in the text.

\paragraph{What has not been done.}
The results adopted from the referee of version~1 were proved by that referee and checked by claude-ab, and they have had no further independent reading. Version~4, including its new results, was read by one further independent referee, a Claude instance that wrote none of the text; its findings were verified and applied. The parts of the workbench marked \textsf{located} were not checked, and §\ref{sec:open} lists what else was left out.

\paragraph{Where everything is.}
This paper, its source and the scripts it rests on are on the branch \texttt{claude/claude-ab-grind-20260925} of \texttt{KokunoYumeto/zeta-function-research-reader}: the paper in \texttt{contrib/claude-ab/rewrite/papers/erdos\_straus/}, the scripts in \texttt{contrib/claude-ab/rewrite/checks/es/}, and the condensed Markdown record of the same material in \texttt{contrib/claude-ab/rewrite/condensed/}. The earlier versions, the audit notes \note{21\_}, \note{43\_}, \note{43a\_}, \note{47\_}, \note{47a\_} and \note{50\_}, and the folders \note{checks/es\_reader/} and \note{checks/es50/} with the referee reports, are in the branch's history at commit \texttt{a90d9e5b}, in \texttt{contrib/claude-ab/grind\_20260925/}.
